\subsection{代数的线性表示}

定义 1. —— 设 M 是一个 K-模。代数 A 在 M 中的一个线性表示是一个 K-同态 \(\pi\)，它从 A 到代数 \(\operatorname{End}_{\mathrm{K}}(\mathrm{M})\)。

我们也说偶 \((\mathrm{M},\pi)\) 是代数 A 的一个线性表示。当 M 是自由 K-模时，M 的维数称为 \(\pi\) 的次数（或维数）。

设 \(\pi\) 是 A 在 M 中的一个线性表示。M 上的加法法则与外法则 \((a,x)\mapsto\pi(a)(x)\) 在 M 上定义了一个左 A-模结构，我们把它记作 \(M_{\pi}\)。我们说 \(M_{\pi}\) 是表示 \(\pi\) 的模。M 上的 K-模结构由 \(M_{\pi}\) 上的 A-模结构通过限制标量得到。

反之，设 E 是一个左 A-模。设 M 是由 E 经标量从 A 到 K 的限制得到的 K-模，\(\pi\) 是同态 \(a \mapsto a_{E}\)，它从 A 到 \(\operatorname{End}_{\mathrm{K}}(\mathrm{M})\)。则 \(\pi\) 是 A 在 M 中的一个表示，且我们有 \(E = M_{\pi}\)。我们说 \(\pi\) 是与 A-模 E 相伴的线性表示。

研究 A-模或 A 的线性表示是一回事。我们将把某些有关模的定义翻译成线性表示的语言。

设 \(\pi\) 是 A 在 M 中的一个线性表示。同态 \(\pi\) 的核是 A 的一个双边理想，它就是 A-模 \(\mathrm{M}_{\pi}\) 的零化子。同态 \(\pi\) 是单射的，当且仅当 A-模 \(\mathrm{M}_{\pi}\) 是忠实的；这时我们说 \(\pi\) 是 A 的一个忠实表示。

设 \((\mathrm{M},\pi)\) 与 \((\mathrm{M}^{\prime},\pi^{\prime})\) 是 A 的线性表示。一个从 \(\mathrm{M}_{\pi}\) 到 \(\mathrm{M}_{\pi^{\prime}}^{\prime}\) 的 A-线性映射是一个 K-线性映射 \(u:\mathrm{M}\to\mathrm{M}^{\prime}\)，它满足关系

\[
\pi^ {\prime} (a) \circ u = u \circ \pi (a)\tag{1}
\]

（对 \(a \in \mathbf{A}\)）。从 \(\mathrm{M}_{\pi}\) 到 \(\mathrm{M}_{\pi'}'\) 的一个同构是 K-模的一个同构 \(u: \mathrm{M} \to \mathrm{M}'\)，它满足条件

\[
\pi^ {\prime} (a) = u \circ \pi (a) \circ u ^ {- 1}\tag{2}
\]

（对 \(a \in A\)）。我们说 \(\pi\) 与 \(\pi'\) 是同构的，如果 A-模 \(M_{\pi}\) 与 \(M'_{\pi'}\) 同构。我们说 \(\pi\) 是 \(\pi'\) 的一个子表示（相应地，商表示），如果 \(M_{\pi}\) 是 \(M'_{\pi'}\) 的一个子模（相应地，商模）。

给定 A 的线性表示的一个族 \((\mathrm{M}_{i},\pi_{i})\)。设 M 是诸 \(M_{i}\) 的直和这个 K-模；对任何 \(a\in A\)，设 \(\pi(a)\) 是 M 的自同态 \((x_{i})_{i\in\mathrm{I}}\mapsto(\pi_{i}(a)(x_{i}))_{i\in\mathrm{I}}\)。则 \(\pi\) 是 A 在 M 中的一个线性表示。A-模 \(M_{\pi}\) 是诸 A-模 \((\mathrm{M}_{i})_{\pi_{i}}\)（\(i\in I\)）的直和；我们说 \(\pi\) 是诸 \(\pi_{i}\) 的直和，并记 \(\pi=\bigoplus\pi_{i}\)。

我们说 A 在 M 中的表示 \(\pi\) 是单的或不可约的，如果 A-模 \(M_{\pi}\) 是单的；我们说 A 在 M 上的表示 \(\pi\) 是半单的或完全可约的，如果 A-模 \(M_{\pi}\) 是半单的。

设 \(\pi\) 是 A 在 M 中的一个线性表示。设 \(\mathrm{M}^*\) 是与 M 对偶的 K-模。映射 \(a\mapsto{}^{t}(\pi (a))\)（它从 \(\mathrm{A}^{\circ}\) 到 \(\mathrm{End}_{\mathrm{K}}(\mathrm{M}^{*})\)）是 \(\pi\) 的转置表示。

设 L 是一个交换 K-代数，\((M,\pi)\) 是 K-代数 A 的一个线性表示。同态 \(\pi_{(L)}:A_{(L)}\to\operatorname{End}_{L}(M_{(L)})\)（与 \(A_{(L)}\) -模在 \(M_{(L)}\) 上的结构相对应）是 L-代数 \(A_{(L)}\) 的一个线性表示。我们说 \(\pi_{(L)}\) 是代数 \(A_{(L)}\) 的、由表示 \(\pi\) 经把标量环 K 扩张到 L 而得到的线性表示。

设 K 是一个体且 L 是一个非零交换 K-代数。

设 \(\pi\) 与 \(\pi'\) 是代数 A 的线性表示。由 VIII, 第 37 页，定理 3，表示 \(\pi\) 与 \(\pi'\) 同构当且仅当 \(\pi_{(L)}\) 与 \(\pi_{(L)}'\) 同构。

设 K 是一个体。设 L 是 K 的一个扩张。考虑格罗滕迪克群 \(R_{\mathrm{K}}(\mathrm{A})\)（相应地，\(R_{\mathrm{L}}(\mathrm{A}_{(\mathrm{L})})\)）——它们分别由在 K 上有限维的 A-模与在 L 上有限维的 \(A_{(L)}\) -模组成。我们已经看到，群同态

\[
u: \mathrm{R} _ {\mathrm{K}} (\mathrm{A}) \longrightarrow \mathrm{R} _ {\mathrm{L}} (\mathrm{A} _ {(\mathrm{L})})
\]

（由标量扩张所定义）是单射的；此外，一个元素 \(\xi \in \mathrm{R}_{\mathrm{K}}(\mathrm{A})\) 是有效的，当且仅当 \(u(\xi)\) 是有效的（VIII, 第 195 页，定理 1）。

